% ED26-SRC-027-501 continuation. Integrator-owned placement in source section3.
% Does not edit or supersede the released core while integration is active.
\subsection{The marked-tree argument at its original hypotheses}
\label{sec:p12-klop1987-marked-tree-proofs}

The following expands the argument in Klop's 1987 report,
\S3, pp.18--22 \parencite{Klop1987CWITutorial}. Its plain objects are
finite, nonempty, rooted trees with natural-number labels and unordered
children. Operator arities are unrestricted in this part of the report.
The stars are temporary control marks, not additional numbers in the
label order. A full tracing proof of the missing acyclicity step is
provided here with the separately attributed argument of
\textcite[\S4.3]{EndrullisKlopOverbeek2021MarkedTrees}; this later proof
dependency is not identified with the 1992 chapter or with a proof
printed in the 1987 tutorial.

\begin{hbdefinition}{def:p12-klop1987-mt-objects}{Unordered tree objects and well-defined marks}
Represent a tree by \(n(t_1,\ldots,t_k)\), with \(n\in\mathbb N\)
and any finite \(k\geq0\), identifying representations by permutations
of siblings at every node. This identification is denoted \(\simeq\)
on representations. Its equivalence classes form \(\mathcal T\).
Allowing each node either its plain label \(n\) or its marked version
\(n^*\) gives \(\mathcal T^*\). Each node has at most one mark.
For a tree with unmarked root \(t=n(\vec u)\), \(t^*=n^*(\vec u)\)
marks that root; it does not recursively mark the children.

The report's operations are putting this root mark; selecting one
child of a marked root; replacing \(n^*(\vec u)\) by
\(m(n^*(\vec u),\ldots,n^*(\vec u))\) when \(m<n\), with any
finite number of copies, including zero; and pushing the mark to a
chosen unmarked-root child \(v\), replacing
\(n^*(v,\vec u)\) by \(n(v^*,\ldots,v^*,\vec u)\), again with
zero or more copies. Thus the last operation removes the parent's mark.
The type of \(v^*\) requires that its root was unmarked; a second mark
on an already marked root is not silently added.
The source reduction is \(\to\), with its stated context closure.
For a proof comparison only, write \(\rightsquigarrow\) for the
same operations closed under \emph{all} finite contexts, including
marked ones. Every source step is a \(\rightsquigarrow\)-step.
This auxiliary enlargement makes no claim that a printed unmarked
context clause and every marked context are the same definition.
\end{hbdefinition}

\begin{hbdefinition}{def:p12-klop1987-mt-embedding}{The report's ancestry-order embedding}
For nodes \(p,q\) of a rooted tree, let \(p\leq_{\rm anc}q\)
mean that \(p\) is an ancestor of \(q\), including equality.
Write \(s\sqsubseteq t\) when there is an injective node map
\(\phi\) satisfying
\[
 p\leq_{\rm anc}q\quad\Longleftrightarrow\quad
 \phi(p)\leq_{\rm anc}\phi(q),\qquad
 \operatorname{label}(p)\leq\operatorname{label}(\phi(p)).
\]
The root may map to a proper descendant of the target root.
The equivalence on the first line makes the map an isomorphism onto
its induced ancestry-ordered image, as specified in the report.
It need not preserve the least common ancestor of two incomparable
nodes. This exact distinction is material to its Figure~10.

For the well-quasi-order proof, temporarily give children a planar
order and use a stronger relation \(\preccurlyeq_{\rm pl}\).
It permits diving into a target child; or coupling roots with
nondecreasing labels and embedding the source child list into a
subsequence of the target child list, using the same relation on
each paired child. Coupling places different source children in
different immediate target branches. Every such planar embedding,
after forgetting sibling order, is an ancestry-order embedding above.
\end{hbdefinition}

\begin{hblemma}{thm:p12-klop1987-mt-embedding}{proved=annex:p12-klop1987-mt-embedding}{Well-defined embedding conventions and their distinct strength}
The source relation \(\sqsubseteq\) is a reflexive, transitive,
antisymmetric relation on unordered labelled-tree classes. The planar
relation is a preorder, and every planar embedding induces a source
embedding. Requiring least-common-ancestor preservation is stronger
than the ancestry-order convention, not an alternative notation for it.
\end{hblemma}

\begin{hbexample}{def:p12-klop1987-mt-figures8-9}{The actual marked trees and the copying example}
Figure~8's plain tree is \(3(5,7(9),8(0(1,5)))\), and its marked
version is \(3^*(5,7(9^*),8^*(0(1,5)))\). The marks are separate
node bits. For Figure~9 put \(a=5(6,7(8))\). Its construction is
\begin{align*}
 a&\to a^*\to4(a^*,a^*,a^*)\\
 &\to^*4(6,5(6,7^*(8)),4(a^*,a^*))\\
 &\to^*4(6,5(6,8),4(6,7(8))).
\end{align*}
The middle segment selects 6 in the first copy, pushes once into 7(8)
in the second, and copies the third twice under 4. The last segment
selects 8 from 7*, then 6 and 7(8) from the last marked 5-copies.
The displayed finite construction ends unmarked. Repeatedly copying
a still marked 5-subtree instead gives an infinite auxiliary reduction.
\end{hbexample}

\begin{hbexample}{def:p12-klop1987-mt-figure10}{The actual 1987 embedding example and the later variant}
The report has
\[
 2(9,7(0,4))\sqsubseteq1(3(8(0(5,1)),9,5(9)),2).
\]
Map its root 2 to the 3-node; its 9-child to that node's 9-child;
its 7-child to the 8-node; and the 0- and 4-children to the 1- and 5-leaves
below that 8-node, respectively. Each label inequality holds, and
ancestry is preserved and reflected. The last two images have the 0-node
as their least common ancestor, not the image 8 of their source parent 7.
Thus this example must not be used to assert meet preservation.
The later primary paper uses instead an 8 in place of that internal 0
and an explicitly meet-preserving embedding
\parencite[Example4.3]{EndrullisKlopOverbeek2021MarkedTrees}.
The stronger planar relation below is used to prove a well-quasi-order
for the report's weaker relation, without collapsing these variants.
\end{hbexample}

\begin{hbdefinition}{def:p12-klop1987-mt-wqo}{The required quasi-order property}
A quasi-order is reflexive and transitive. It is a well-quasi-order
(wqo) if every infinite sequence has indices \(i<j\) with its
\(i\)-th member below its \(j\)-th member. A bad sequence lacks
such a pair. Finite words are ordered by increasing subsequence
embedding, allowing deletion of letters and replacement by smaller
letters. Empty words are included.
All alphabets used in the following constructions have a countable
coding: natural labels, finite trees over those labels, and their
finite lists. Minimal bad sequences are chosen by length or number
of nodes, with a fixed enumeration breaking ties; this is a classical
existence argument, not an algorithm deciding whether a bad extension
exists.
\end{hbdefinition}

\begin{hblemma}{thm:p12-klop1987-mt-words}{proved=annex:p12-klop1987-mt-words}{The word and product dependencies}
Every infinite sequence in a wqo has an infinite nondecreasing
subsequence. Finite words over a countably coded wqo form a wqo under
subsequence embedding, and the product of two wqos is a wqo under
componentwise comparison. The natural numbers with their usual order
satisfy these hypotheses.
\end{hblemma}

\begin{hbtheorem}{thm:p12-klop1987-mt-kruskal}{proved=annex:p12-klop1987-mt-kruskal}{The full labelled-tree wqo needed in the report}
Finite planar natural-number-labelled trees are wqo under
\(\preccurlyeq_{\rm pl}\). Consequently the report's unordered
trees are wqo under \(\sqsubseteq\). Labels, depths and branching
degrees need not have one common finite bound. This proves the
natural-number version of the dependency invoked at Theorem~3.8,
not merely a fixed finite-ranked-alphabet special case.
\end{hbtheorem}

\begin{hblemma}{thm:p12-klop1987-mt-embedding-reduction}{proved=annex:p12-klop1987-mt-embedding-reduction}{Embedding is simulated in the report's marker relation}
If \(s\sqsubseteq t\), then \(t\to^*s\). If the tree classes
are different, the reduction is positive. The identical case uses
zero steps; the unrestricted assertion \(t\to^+t\) is not inferred
from reflexive embedding. The simulation handles distinct source
branches whose image roots share an initial target stem.
\end{hblemma}

\begin{hblemma}{thm:p12-klop1987-mt-tracing}{proved=annex:p12-klop1987-mt-tracing}{The occurrence-level marker mechanism}
In the all-context auxiliary game, each output node has one input
ancestor. Its numerical label does not exceed that ancestor's label;
a new lower root introduced by copying is a strict decrease on that
lineage. A tracked subtree has a projected reduction using at most
one step per original step. Finite sibling permutations can be
postponed, so a cycle on tree classes would give a marker reduction
on planar representations ending in a permutation of its input.
\end{hblemma}

\begin{hbtheorem}{thm:p12-klop1987-mt-acyclic}{proved=annex:p12-klop1987-mt-acyclic}{Acyclicity on plain trees, including the zero-copy cases}
There is no positive \(\rightsquigarrow\)-reduction between an
unmarked tree and itself. Hence the report's source relation has
no such cycle. The assertion is not acyclicity or termination of
the whole marked auxiliary system.
\end{hbtheorem}

\begin{hbtheorem}{thm:p12-klop1987-mt-well-founded}{proved=annex:p12-klop1987-mt-well-founded}{The report's unmarked strict order is well-founded}
The relation \(\to^+\) restricted to pairs in \(\mathcal T\)
is a strict, transitive, well-founded order. In particular there is
no infinite sequence of unmarked trees joined successively by
positive finite marker reductions, as asserted in Theorem~3.5.
The proof uses the established wqo, embedding simulation and
acyclicity; none is replaced by its name.
\end{hbtheorem}

The marked universe is genuinely different. For example,
\(1^*()\to0(1^*())\to0(0(1^*()))\to\cdots\)
uses repeated copying with one copy and never becomes entirely
unmarked. It witnesses failure of SN on \(\mathcal T^*\), exactly
the distinction made before the report's Theorem~3.5.

\begin{hbtheorem}{thm:p12-klop1987-mt-boolean}{proved=annex:p12-klop1987-mt-boolean}{All five source Boolean-formula rewrite comparisons}
The five rules of source Table~9 terminate: double-negation removal,
both De Morgan transformations, and both distributive transformations
of conjunction over disjunction. Their positive marker simulations
use \(\lor\mapsto1\), \(\land\mapsto2\), \(\neg\mapsto3\).
This is termination of these syntactic rules, not confluence or
completeness of a logical proof calculus.
\end{hbtheorem}

\begin{hbtheorem}{thm:p12-klop1987-mt-indexed}{proved=annex:p12-klop1987-mt-indexed}{The norm-indexed source pair, with the context dependency checked}
For the two rules \(g(x,y)\to h(x,y)\) and
\(h(f(x),y)\to f(g(x,y))\), the source's norm-indexed construction
proves termination. Annotate a binary \(g(s,r)\) or \(h(s,r)\)
by \(n=|s|+|r|\), with bars counting nodes, and take labels
\(g_n=2n\), \(h_n=2n-1\), \(f=0\).
Then \(g_n>h_n>f\) and \(h_{n+1}>g_n\).
Both original rules preserve node count, so ancestor annotations
stay unchanged when a rule is used in an arbitrary context.
This establishes the concrete two-rule construction; it does not
silently extend node-count preservation to the full process-algebra
rules that also duplicate arguments under a sum.
\end{hbtheorem}

\begin{hbcorollary}{thm:p12-klop1987-mt-obstruction}{proved=annex:p12-klop1987-mt-obstruction}{The source's embedded-left-side obstruction}
An ordinary rule whose left side embeds in its right side cannot have
a positive descending simulation in the same acyclic marker order.
The source example \(f(s(x))\to g(s(x),f(p(s(x))))\) has that
obstruction whenever each operator's label is unchanged between its
two occurrences.
\end{hbcorollary}

\begin{hbdefinition}{def:p12-klop1987-mt-ordered-extension}{The separately enlarged ordered marker relation}
Remark~3.11(iii) changes the objects to ordered trees with fixed
operator arities. Copying to a lower label must give that lower
operator its prescribed number of arguments. Pushing changes one
chosen argument \(t_i\) to \(t_i^*\), without deleting or copying
that argument. For \(n\) of arity \(k\geq1\), the added
left-argument operation is
\[
 n^*(t_1,\ldots,t_k)\longrightarrow
 n(t_1^*,n^*(t_1,\ldots,t_k),\ldots,n^*(t_1,\ldots,t_k)),
\]
with exactly \(k-1\) full copies after the first argument.
The marker on \(t_1\) is defined at an unmarked root as above.
The report gives the Ackermann comparison
\begin{align*}
 A(S(x),S(y))&\longrightarrow A^*(S(x),S(y))\\
 &\longrightarrow A(S^*(x),A^*(S(x),S(y)))\\
 &\longrightarrow A(x,A^*(S(x),S(y)))\\
 &\longrightarrow A(x,A(S(x),S^*(y)))\\
 &\longrightarrow A(x,A(S(x),y)).
\end{align*}
Its source rules are \(A(0,x)\to S(x)\),
\(A(S(x),0)\to A(x,S(0))\), and
\(A(S(x),S(y))\to A(x,A(S(x),y))\).
This displayed derivation belongs to the enlarged ordered relation,
not to the unranked commutative relation proved well-founded above.
\end{hbdefinition}

\begin{proofboundary}{def:p12-klop1987-mt-boundaries}{Separate mechanisms and separate source standing}
The main proof concerns the unranked natural-labelled source tree
order. Adding the left-argument operation changes the reduction
relation, so inclusion in the proved four-operation game is unavailable.
The global well-foundedness of that enlarged relation is instead
established by the separate lower-set interpretation in
Theorem~\ref{thm:p12-klop1987-om-well-founded}; its strict conclusion
concerns plain endpoints, not the whole marked universe. An Ackermann
marker comparison alone would not supply that argument. Likewise, the displayed
rule whose left side embeds in its right side requires a different
method: \(f(s(x))\to g(s(x),f(p(s(x))))\) cannot strictly descend
in an acyclic order also containing reverse embedding.
The report refers that latter extension to Puel(1986), and refers
the ordered lexicographic extension to Kamin--Levy through Dershowitz.
The ordered-extension proof does not discharge the distinct general
Puel extension. The actual three ordinary PA rules instead have the
separate contextual max/sum and retained-budget proofs in
Theorems~\ref{thm:p12-klop1987-pa-termination} and
\ref{thm:p12-klop1987-pa-budget-lift}, at their declared three-rule
and seven-rule scopes. The larger ACP communication/silent-action
variants and Puel's general unavoidable-set construction remain outside
these proofs; neither is established by the isolated example's
active-origin argument.
\end{proofboundary}
